\documentclass[../../main.tex]{subfiles}
\begin{document}

\section{Seifert-van Kampen 定理}\label{section:ch11-Seifert-van-Kampen定理}

我们现在回到决定空间 $X$ 的基本群的问题，这里空间 $X$ 可以表示成两个开集 $U$ 和 $V$ 之并，并且它们的交是道路连通的。我们已在第 59 节中证明了：如果 $x_0\in U\cap V$，两个群 $\pi_1(U,x_0)$ 和 $\pi_1(V,x_0)$ 在由内射诱导的同态下在 $\pi_1(X,x_0)$ 中的像生成 $\pi_1(X,x_0)$。本节我们将证明 $\pi_1(X,x_0)$ 完全由上述两个群、群 $\pi_1(U\cap V,x_0)$ 以及由内射诱导的群的同态决定。这是关于基本群的一个基本结论。应用这一结论，我们便能够计算包括 2 维紧致流形在内的许多空间的基本群。

\par
\begin{figure}[htbp]
\centering
\includegraphics[width=0.5\textwidth]{image/fig-70-diag1.png}
\caption{}\label{figure:ch11-fig-70-diag1}
\end{figure}

\begin{theorem}[Seifert-van Kampen 定理]\label{theorem:ch11-seifert-van-kampen}
设 $X=U\cup V$，其中 $U$ 和 $V$ 为 $X$ 中的开集。假设 $U$，$V$ 以及 $U\cap V$ 都是道路连通的，$x_0\in U\cap V$。设 $H$ 是一个群，并且
\begin{align*}
\phi_1:\pi_1(U,x_0)\to H\quad\text{和}\quad\phi_2:\pi_1(V,x_0)\to H
\end{align*}
是两个同态。设 $i_1$，$i_2$，$j_1$，$j_2$ 是依照下图所示内射诱导的同态。
如果 $\phi_1\circ i_1=\phi_2\circ i_2$，则存在一个唯一同态 $\Phi:\pi_1(X,x_0)\to H$ 使得 $\Phi\circ j_1=\phi_1$ 和 $\Phi\circ j_2=\phi_2$ 成立。
\end{theorem}

这个定理告诉我们：如果 $\phi_1$ 和 $\phi_2$ 是“在 $U\cap V$ 上相容”的任何两个同态，则它们诱导出从 $\pi_1(X,x_0)$ 到 $H$ 的一个同态。

\par
\begin{figure}[htbp]
\centering
\includegraphics[width=0.5\textwidth]{image/fig-70-1.png}
\caption{}\label{figure:ch11-fig-70-1}
\end{figure}

\par
\begin{figure}[htbp]
\centering
\includegraphics[width=0.7\textwidth]{image/fig-70-2.png}
\caption{}\label{figure:ch11-fig-70-2}
\end{figure}

\begin{proof}
唯一性的证明是简单的。定理 59.1 告诉我们：$\pi_1(X,x_0)$ 是由 $j_1$ 和 $j_2$ 的像生成的。$\Phi$ 在生成元 $j_1(g_1)$ 上的值必为 $\phi_1(g_1)$，它在生成元 $j_2(g_2)$ 上的值必为 $\phi_2(g_2)$。因此 $\Phi$ 完全由 $\phi_1$ 和 $\phi_2$ 所决定。而证明 $\Phi$ 的存在性则不那么容易。

为方便起见，我们引入以下记号：对于 $X$ 中的一条道路 $f$，用 $[f]$ 表示 $f$ 在 $X$ 中的道路同伦类。如果 $f$ 位于 $U$ 中，用 $[f]_U$ 表示 $f$ 在 $U$ 中的道路同伦类。$[f]_V$ 及 $[f]_{U\cap V}$ 等记号的定义是类似的。

\textbf{第一步。} 我们首先定义一个集合间的映射 $\rho$，它将 $U$ 或者 $V$ 中以 $x_0$ 为基点的回路 $f$ 映为群 $H$ 中的一个元素。规定
\begin{align*}
\rho(f)=\phi_1([f]_U)\quad\text{如果 }f\text{ 在 }U\text{ 中},\\
\rho(f)=\phi_2([f]_V)\quad\text{如果 }f\text{ 在 }V\text{ 中}.
\end{align*}
这时 $\rho$ 的定义是确切的，这是因为若 $f$ 既在 $U$ 中又在 $V$ 中，则有
\begin{align*}
\phi_1([f]_U)=\phi_1i_1([f]_{U\cap V})\quad\text{和}\quad\phi_2([f]_V)=\phi_2i_2([f]_{U\cap V}),
\end{align*}
根据假设 $H$ 中的这两个值相等。集合间的映射 $\rho$ 满足以下两个条件：
\begin{enumerate}[(1)]
\item 如果 $[f]_U=[g]_U$ 或者 $[f]_V=[g]_V$，则 $\rho(f)=\rho(g)$。
\item 如果 $f$ 和 $g$ 都在 $U$ 中或者都在 $V$ 中，则 $\rho(f*g)=\rho(f)\cdot\rho(g)$。
\end{enumerate}
前者成立是根据定义，后者成立是因为 $\phi_1$ 和 $\phi_2$ 都是同态。

\textbf{第二步。} 进一步将 $\rho$ 扩充为 $\sigma$，它将每一条 $U$ 或 $V$ 中的道路 $f$ 映为 $H$ 中的某一个元素，要求 $\sigma$ 满足第一步中的条件 (1)，并且当 $f*g$ 有定义时，$\sigma$ 满足条件 (2)。

首先，对于每一个 $x\in X$，依如下方式选取一条从 $x_0$ 到 $x$ 的道路 $\alpha_x$：如果 $x=x_0$，取 $\alpha_x$ 为 $x_0$ 点处的常值道路。如果 $x\in U\cap V$，取 $\alpha_x$ 为 $U\cap V$ 中的某一条道路。如果 $x$ 在 $U$ 或 $V$ 中但不在 $U\cap V$ 中，分别取 $\alpha_x$ 为 $U$ 或者 $V$ 中的某一条道路。

然后，对任何 $U$ 或 $V$ 中的道路 $f$，按照
\begin{align*}
L(f)=\alpha_x*(f*\bar{\alpha}_y)
\end{align*}
分别定义 $U$ 或 $V$ 中以 $x_0$ 为基点的回路 $L(f)$，其中 $x$ 为 $f$ 的起点，$y$ 为 $f$ 的终点。见图 70.1。最后，定义
\begin{align*}
\sigma(f)=\rho(L(f)).
\end{align*}


我们先来证明 $\sigma$ 的确是 $\rho$ 的一个扩充。如果 $f$ 是位于 $U$ 或 $V$ 中的一条以 $x_0$ 为基点的回路，则由于 $\alpha_{x_0}$ 是常值道路，从而
\begin{align*}
L(f)=e_{x_0}*(f*e_{x_0}).
\end{align*}
那么，$L(f)$ 在 $U$ 或 $V$ 中道路同伦于 $f$。由于 $\rho$ 满足条件 (1)，所以 $\rho(L(f))=\rho(f)$。因此 $\sigma(f)=\rho(f)$。

为了验证条件 (1)，设 $f$ 和 $g$ 是位于 $U$ 或者 $V$ 中彼此道路同伦的道路。那么，回路 $L(f)$ 和 $L(g)$ 也是在 $U$ 或者 $V$ 中彼此道路同伦的，从而 $\rho$ 满足条件 (1)。为了验证条件 (2)，设 $f$ 和 $g$ 为 $U$ 或 $V$ 中满足 $f(1)=g(0)$ 的任何两条道路，则对于适当的点 $x$，$y$ 以及 $z$，有
\begin{align*}
L(f)*L(g)=(\alpha_x*(f*\bar{\alpha}_y))*(\alpha_y*(g*\bar{\alpha}_z)),
\end{align*}
这条回路在 $U$ 或 $V$ 中道路同伦于 $L(f*g)$。对于 $\rho$ 应用条件 (1) 和 (2)，则有
\begin{align*}
\rho(L(f*g))=\rho(L(f)*L(g))=\rho(L(f))\cdot\rho(L(g)).
\end{align*}
因此 $\sigma(f*g)=\sigma(f)\cdot\sigma(g)$。

\textbf{第三步。} 最后，我们将 $\sigma$ 扩充为集合间的映射 $\tau$，使得 $\tau$ 将 $X$ 中的任意一条道路 $f$ 映为 $H$ 中的一个元素，要求它满足以下条件：
\begin{enumerate}[(1)]
\item 如果 $[f]=[g]$，则 $\tau(f)=\tau(g)$。
\item 如果 $f*g$ 有定义，则 $\tau(f*g)=\tau(f)\cdot\tau(g)$。
\end{enumerate}

对于给定的 $f$，选取 $[0,1]$ 区间的一个分划 $s_0<\cdots<s_n$，使得 $f$ 将每一个子区间 $[s_{i-1},s_i]$ 映到 $U$ 或 $V$ 中。令 $f_i$ 表示 $[0,1]$ 到 $[s_{i-1},s_i]$ 的正线性映射与 $f\mid_{[s_{i-1},s_i]}$ 的复合。这时 $f_i$ 是位于 $U$ 或 $V$ 中的道路，并且
\begin{align*}
[f]=[f_1]*\cdots*[f_n].
\end{align*}
如果 $\tau$ 是 $\sigma$ 的一个扩充并且满足条件 (1) 和 (2)，则必有
\begin{align}
\tau(f)=\sigma(f_1)\cdot\sigma(f_2)\cdots\sigma(f_n).\label{align:ch11-svk-tau-def}
\end{align}
于是我们将用此等式作为 $\tau$ 的定义。

我们将证明这个定义与分划的选取无关。这只要证明：若将某一个点 $p$ 插入到分划中，则 $\tau(f)$ 的值保持不变。设指标 $i$ 满足条件 $s_{i-1}<p<s_i$，当对于新分划计算 $\tau(f)$ 时，公式 \eqref{align:ch11-svk-tau-def} 的变化仅仅是 $\sigma(f_i)$ 不出现，取而代之的是 $\sigma(f_i')\cdot\sigma(f_i'')$，其中 $f_i'$ 和 $f_i''$ 分别为 $[0,1]$ 到 $[s_{i-1},p]$ 和 $[p,s_i]$ 的正线性映射与 $f$ 在相应区间上的限制的复合。而 $f_i$ 在 $U$ 或 $V$ 中与 $f_i'*f_i''$ 道路同伦，由于 $\sigma$ 满足条件 (1) 及 (2)，我们有 $\sigma(f_i)=\sigma(f_i')\cdot\sigma(f_i'')$。因此，$\tau$ 的定义是确切的。

因而，$\tau$ 是 $\sigma$ 的一个扩充。因为若 $f$ 本身在 $U$ 或 $V$ 中，我们可以选取 $[0,1]$ 上的平凡分划来定义 $\tau(f)$，从而根据定义，$\tau(f)=\sigma(f)$。

\textbf{第四步。} 我们证明条件 (1) 对于集合间的映射 $\tau$ 成立。这部分的证明需特别小心。

我们先就一个特殊情形验证这个条件。设 $f$ 和 $g$ 为 $X$ 中由 $x$ 到 $y$ 的两条道路，$F$ 为二者间的同伦。此外，我们还假设存在 $[0,1]$ 的分划 $s_0,\cdots,s_n$，使得 $F$ 将每一个矩形 $R_i=[s_{i-1},s_i]\times I$ 映入 $U$ 与 $V$ 之一。我们在这个附加的假设下来证明 $\tau(f)=\tau(g)$。

对于给定的 $i$，考虑由 $[0,1]$ 到 $[s_{i-1},s_i]$ 的正线性映射与 $f$ 或 $g$ 的复合，将它们分别记作 $f_i$ 和 $g_i$。$F$ 在矩形 $R_i$ 上的限制给出的是 $f_i$ 与 $g_i$ 间在 $U$ 或 $V$ 中的一个同伦，然而 $F$ 未必是一个道路同伦，因为在同伦过程中道路的端点可能移动。我们来考察同伦过程中道路端点的轨迹。定义 $\beta_i$ 为道路 $\beta_i(t)=F(s_i,t)$，则 $\beta_i$ 为 $X$ 中从 $f(s_i)$ 到 $g(s_i)$ 的一条道路。道路 $\beta_0$ 和 $\beta_n$ 分别为点 $x$ 和 $y$ 处的常值道路。见图 70.2。我们证明：对于每一个 $i$，
\begin{align*}
f_i*\beta_i\simeq_p\beta_{i-1}*g_i
\end{align*}
是在 $U$ 或 $V$ 中的道路同伦。


在矩形 $R_i$ 中，我们取沿 $R_i$ 的下边线和右边线的由 $s_{i-1}\times 0$ 经由 $s_i\times 0$ 到 $s_i\times 1$ 的那条折线道路，这条道路在 $F$ 的作用下便得到了道路 $f_i*\beta_i$。类似地，取沿 $R_i$ 的左边线和上边线的那条折线道路，这条道路在 $F$ 的作用下便得到了道路 $\beta_{i-1}*g_i$。由于 $R_i$ 是凸集，存在一个在 $R_i$ 中的两条折线间的道路同伦，因此由 $F$ 给出 $f_i*\beta_i$ 与 $\beta_{i-1}*g_i$ 之间的一个在 $U$ 或 $V$ 中的道路同伦。这便是我们所需的道路同伦。

由于 $\sigma$ 满足条件 (1) 和 (2)，我们有
\begin{align*}
\sigma(f_i)\cdot\sigma(\beta_i)=\sigma(\beta_{i-1})\cdot\sigma(g_i),
\end{align*}
因此
\begin{align}
\sigma(f_i)=\sigma(\beta_{i-1})\cdot\sigma(g_i)\cdot\sigma(\beta_i)^{-1}.\label{align:ch11-svk-sigma-relation}
\end{align}
类似地，由于 $\beta_0$ 和 $\beta_n$ 为常值道路，我们有 $\sigma(\beta_0)=\sigma(\beta_n)=1$。（因为 $\beta_0*\beta_0=\beta_0$ 蕴涵着 $\sigma(\beta_0)\cdot\sigma(\beta_0)=\sigma(\beta_0)$。）

推导如下：
\begin{align*}
\tau(f)=\sigma(f_1)\cdot\sigma(f_2)\cdots\sigma(f_n).
\end{align*}
对此式应用 \eqref{align:ch11-svk-sigma-relation} 并化简，我们有
\begin{align*}
\tau(f)&=\sigma(g_1)\cdot\sigma(g_2)\cdots\sigma(g_n)\\
&=\tau(g).
\end{align*}
从而，我们就这一特别情形证明了条件 (1)。

下面我们就一般情形证明条件 (1)。任意给定道路 $f$ 和 $g$ 以及二者间的一个同伦 $F$，选取 $[0,1]$ 的分划 $s_0,\cdots,s_n$ 和 $t_0,\cdots,t_m$ 使得 $F$ 将每一个子矩形 $[s_{i-1},s_i]\times[t_{j-1},t_j]$ 映入 $U$ 或 $V$ 中。若 $f_j$ 为道路 $f_j(s)=F(s,t_j)$，则 $f_0=f$ 且 $f_m=g$。由于道路 $f_{j-1}$ 和 $f_j$ 满足前述特殊情形的条件，从而对于每一个 $j$ 有 $\tau(f_{j-1})=\tau(f_j)$，所以有 $\tau(f)=\tau(g)$。

\textbf{第五步。} 我们证明集合之间的映射 $\tau$ 满足条件 (2)。给定 $X$ 中的一条道路 $f*g$，选取 $[0,1]$ 的一个包含点 $1/2$ 作为分点的分划 $s_0<\cdots<s_n$，使得 $f*g$ 将每一子区间映射到 $U$ 或 $V$ 中。设指标 $k$ 满足条件 $s_k=1/2$。

当 $i=1,\cdots,k$ 时，由 $[0,1]$ 到 $[s_{i-1},s_i]$ 的正线性映射复合 $f*g$ 所得到的道路就是由 $[0,1]$ 到 $[2s_{i-1},2s_i]$ 的正线性映射复合 $f$ 所得到的道路，将此道路记为 $f_i$。类似地，当 $i=k+1,\cdots,n$ 时，由 $[0,1]$ 到 $[s_{i-1},s_i]$ 的正线性映射复合 $f*g$ 所得到的道路就是由 $[0,1]$ 到 $[2s_{i-1}-1,2s_i-1]$ 的正线性映射复合 $g$ 所得到的道路，将此道路记为 $g_{i-k}$。将分划 $s_0,\cdots,s_n$ 应用于道路 $f*g$ 的定义域，我们有
\begin{align*}
\tau(f*g)=\sigma(f_1)\cdots\sigma(f_k)\cdot\sigma(g_1)\cdots\sigma(g_{n-k}).
\end{align*}
将分划 $2s_0,\cdots,2s_k$ 应用于道路 $f$，我们有
\begin{align*}
\tau(f)=\sigma(f_1)\cdots\sigma(f_k).
\end{align*}
将分划 $2s_k-1,\cdots,2s_n-1$ 应用于道路 $g$，我们有
\begin{align*}
\tau(g)=\sigma(g_1)\cdots\sigma(g_{n-k}).
\end{align*}
因此 (2) 显然成立。

\textbf{第六步。} 定理的证明。对于 $X$ 中以 $x_0$ 为基点的每一条回路 $f$，我们定义
\begin{align*}
\Phi([f])=\tau(f).
\end{align*}
条件 (1) 和 (2) 说明 $\Phi$ 为定义确切的同态。

我们来证明 $\Phi\circ j_1=\phi_1$。如果 $f$ 是 $U$ 中的一条回路，则
\begin{align*}
\Phi(j_1([f]_U))&=\Phi([f])\\
&=\tau(f)\\
&=\rho(f)=\phi_1([f]_U).
\end{align*}
类似地，$\Phi\circ j_2=\phi_2$。
\end{proof}

上述定理是 Seifert-van Kampen 定理的现代形式。以下我们再来给出它的经典形式，它涉及两个群的自由积。前面我们已经提到过，当 $G$ 是外自由积 $G=G_1*G_2$ 时，为简化记号，经常将 $G_1$ 和 $G_2$ 视为 $G$ 的子群。

\begin{theorem}[Seifert-van Kampen 定理的经典形式]\label{theorem:ch11-seifert-van-kampen-classical}
假设如前一定理。设
\begin{align*}
j:\pi_1(U,x_0)*\pi_1(V,x_0)\to\pi_1(X,x_0)
\end{align*}
是自由积的同态，并且是由内射诱导的同态 $j_1$ 和 $j_2$ 的扩充。则 $j$ 是一个满射，并且它的核为自由积的最小正规子群 $N$，它包含了所有形如
\begin{align*}
(i_1(g)^{-1},i_2(g))
\end{align*}
的字表示的元素，其中 $g\in\pi_1(U\cap V,x_0)$。

换言之，$j$ 的核是由自由积中所有形如 $i_1(g)^{-1}i_2(g)$ 的元素以及它们的共轭元所生成的群。
\end{theorem}

\par
\begin{figure}[htbp]
\centering
\includegraphics[width=0.5\textwidth]{image/fig-70-diag2.png}
\caption{}\label{figure:ch11-fig-70-diag2}
\end{figure}

\begin{proof}
由于 $\pi_1(X,x_0)$ 是由 $j_1$ 和 $j_2$ 的像所生成的，因此 $j$ 为满射。

我们证明 $N\subset\ker j$。由于 $\ker j$ 是正规的，我们只要证明对任何 $g\in\pi_1(U\cap V,x_0)$，$i_1(g)^{-1}i_2(g)$ 属于 $\ker j$。如果 $i:U\cap V\to X$ 是内射，则
\begin{align*}
ji_1(g)=j_1i_1(g)=i_*(g)=j_2i_2(g)=ji_2(g).
\end{align*}
于是 $i_1(g)^{-1}i_2(g)$ 属于 $j$ 的核。

因此 $j$ 诱导一个满同态
\begin{align*}
k:\pi_1(U,x_0)*\pi_1(V,x_0)/N\to\pi_1(X,x_0).
\end{align*}
我们证明 $k$ 是一个单射，由此得到 $N=\ker j$。这只要证明 $k$ 有一个左逆。

令 $H$ 表示群 $\pi_1(U,x_0)*\pi_1(V,x_0)/N$。又设 $\phi_1:\pi_1(U,x_0)\to H$ 是 $\pi_1(U,x_0)$ 到自由积的内射与自由积到它相对于 $N$ 的商群的投射两者的复合。类似地，定义 $\phi_2:\pi_1(V,x_0)\to H$。考虑下图表：

易见，$\phi_1\circ i_1=\phi_2\circ i_2$。因为若 $g\in\pi_1(U\cap V,x_0)$，则 $\phi_1(i_1(g))$ 为 $H$ 中的陪集 $i_1(g)N$，而 $\phi_2(i_2(g))$ 为陪集 $i_2(g)N$。由于 $i_1(g)^{-1}i_2(g)\in N$，所以两个陪集相等。

根据定理 70.1 推出：存在一个同态 $\Phi:\pi_1(X,x_0)\to H$ 使得 $\Phi\circ j_1=\phi_1$ 及 $\Phi\circ j_2=\phi_2$。我们证明 $\Phi$ 是 $k$ 的左逆。这只要证明 $\Phi\circ k$ 作用在 $H$ 的每一个生成元上恒等，也就是说，对每一个属于 $\pi_1(U,x_0)$ 或 $\pi_1(V,x_0)$ 的 $g$，$\Phi\circ k$ 作用在形如 $gN$ 的陪集上等于其自身。因为当 $g\in\pi_1(U,x_0)$ 时，我们有
\begin{align*}
k(gN)=j(g)=j_1(g),
\end{align*}
由此得到
\begin{align*}
\Phi(k(gN))=\Phi(j_1(g))=\phi_1(g)=gN,
\end{align*}
这便是我们要证明的。当 $g\in\pi_1(V,x_0)$ 时情况是类似的。
\end{proof}

\begin{corollary}\label{corollary:ch11-svk-simply-connected-intersection}
假设同 Seifert-van Kampen 定理。如果 $U\cap V$ 是单连通的，则存在一个同构
\begin{align*}
k:\pi_1(U,x_0)*\pi_1(V,x_0)\to\pi_1(X,x_0).
\end{align*}
\end{corollary}

\begin{corollary}\label{corollary:ch11-svk-simply-connected-v}
假设同 Seifert-van Kampen 定理。如果 $V$ 是单连通的，则存在一个同构
\begin{align*}
k:\pi_1(U,x_0)/N\to\pi_1(X,x_0),
\end{align*}
其中 $N$ 为包含同态
\begin{align*}
i_1:\pi_1(U\cap V,x_0)\to\pi_1(U,x_0)
\end{align*}
的像的 $\pi_1(U,x_0)$ 的最小正规子群。
\end{corollary}

\par
\begin{figure}[htbp]
\centering
\includegraphics[width=0.3\textwidth]{image/fig-70-3.png}
\caption{}\label{figure:ch11-fig-70-3}
\end{figure}

\begin{example}\label{example:ch11-theta-space}
设 $X$ 是一个 $\theta$ 空间，则 $X$ 是可表示成三条弧 $A$，$B$ 和 $C$ 之并的一个 Hausdorff 空间，其中每两条弧的交点恰为它们各自的端点 $p$ 和 $q$。前面我们已经证明了 $X$ 的基本群不是阿贝尔群。此处我们证明其基本群为具有两个生成元的自由群。

设 $a$ 和 $b$ 分别为 $A$ 和 $B$ 的内点。将 $X$ 写成两个开集 $U=X-a$ 与 $V=X-b$ 之并。见图 70.3。空间 $U\cap V=X-a-b$ 是单连通的，因为它是可缩的。进而，$U$ 和 $V$ 都以无限循环群为基本群，因为 $U$ 有 $B\cup C$ 的伦型，$V$ 有 $A\cup C$ 的伦型。因此，$X$ 的基本群是两个无限循环群的自由积，也就是说，它是具有两个生成元的自由群。
\end{example}

\subsection*{习题}

在以下习题中，假设条件同 Seifert-van Kampen 定理中的一样。

\end{document}